\chapter{Sorting and Big-O Notation}

\section{Why Sorting Matters}\index{sorting}
Sorting a list is one of the most common things a program does. Phone books,
search results, and leaderboards all depend on items being in order. Python
already sorts for you:

\begin{minted}{python}
scores = [38, 27, 43, 3, 9, 82, 10]
print(sorted(scores))  # [3, 9, 10, 27, 38, 43, 82]
\end{minted}

This chapter shows how sorting can be done by hand, using three simple
algorithms. Then it introduces a way to measure how much work an algorithm
does as the list grows: \index{Big-O notation}\newterm{Big-O notation}.

\section{Selection Sort}\index{selection sort}
Selection sort repeats one idea. Find the smallest item in the unsorted part
of the list, and swap it into the next position at the front.

\begin{minted}{python}
def selection_sort(items):
    items = list(items)
    n = len(items)
    for i in range(n - 1):
        smallest = i
        for j in range(i + 1, n):
            if items[j] < items[smallest]:
                smallest = j
        items[i], items[smallest] = items[smallest], items[i]
    return items

print(selection_sort([5, 2, 9, 1, 5, 6]))  # [1, 2, 5, 5, 6, 9]
\end{minted}

The outer loop fixes one position per pass. The inner loop looks at every
remaining item to find the smallest. Notice that the inner loop always runs to
the end of the list, even if the list is already sorted.

\section{Insertion Sort}\index{insertion sort}
Insertion sort works the way many people sort a hand of playing cards. Take
the next card, and slide it left past every larger card until it fits.

\begin{minted}{python}
def insertion_sort(items):
    items = list(items)
    for i in range(1, len(items)):
        current = items[i]
        j = i - 1
        while j >= 0 and items[j] > current:
            items[j + 1] = items[j]
            j -= 1
        items[j + 1] = current
    return items

print(insertion_sort([5, 2, 9, 1, 5, 6]))  # [1, 2, 5, 5, 6, 9]
\end{minted}

Insertion sort is fast when the list is almost sorted, because the inner loop
stops early. When the list is in reverse order, it does the most work.

\section{Merge Sort}\index{merge sort}
Merge sort uses a different strategy: divide and conquer. Split the list in
half, sort each half, and then merge the two sorted halves into one. A list of
one item is already sorted, which is where the recursion stops.

\begin{minted}{python}
def merge_sort(items):
    if len(items) <= 1:
        return list(items)
    mid = len(items) // 2
    left = merge_sort(items[:mid])
    right = merge_sort(items[mid:])
    merged = []
    i = j = 0
    while i < len(left) and j < len(right):
        if left[i] <= right[j]:
            merged.append(left[i]); i += 1
        else:
            merged.append(right[j]); j += 1
    merged.extend(left[i:])
    merged.extend(right[j:])
    return merged

print(merge_sort([38, 27, 43, 3, 9, 82, 10]))  # [3, 9, 10, 27, 38, 43, 82]
\end{minted}

\section{Big-O Notation}\index{Big-O notation!definition}
We want to compare algorithms without timing them on one computer. Big-O
notation describes how the number of steps grows as the input size $n$ grows.
We count the steps in terms of $n$, and then keep only the term that grows
fastest, dropping constants and smaller terms.

We write $f(n) = O(g(n))$ when $f(n)$ is no larger than a constant multiple of
$g(n)$ for all large enough $n$. Formally, there are constants $c > 0$ and
$n_0$ such that $f(n) \leq c \cdot g(n)$ whenever $n \geq n_0$.

For example, $f(n) = 3n^2 + 100n + 5$ is $O(n^2)$. With $c = 4$ and
$n_0 = 101$, we have $3n^2 + 100n + 5 \leq 4n^2$. The $100n$ term matters for
small $n$, but for large $n$ the $n^2$ term dominates.

Here are the classes you will meet most often, from fastest to slowest:

\begin{center}
\begin{tabular}{|l|l|p{6cm}|}
\hline
Notation & Name & Example \\
\hline
$O(1)$ & constant & reading one item from a list by its index \\
\hline
$O(\log n)$ & logarithmic & binary search on a sorted list \\
\hline
$O(n)$ & linear & checking every item once \\
\hline
$O(n \log n)$ & linearithmic & merge sort, and Python's \texttt{sorted()} in the worst case \\
\hline
$O(n^2)$ & quadratic & selection sort and insertion sort \\
\hline
\end{tabular}
\end{center}

\section{Comparing the Sorts}\index{sorting!comparison}
Selection sort's inner loop runs $n - 1$, then $n - 2$, and so on down to 1
times. That is a total of $\frac{n(n-1)}{2}$ comparisons, which is
$O(n^2)$ in every case, even when the list is already sorted. Insertion sort
can do far fewer comparisons on nearly sorted lists, but its worst case, a
list in reverse order, also takes $\frac{n(n-1)}{2}$ comparisons, so it is also
$O(n^2)$.

Merge sort halves the list at each level of recursion, so there are about
$\log_2 n$ levels. Each level does about $n$ work to merge. That gives
$O(n \log n)$ work in every case.

The table below shows the approximate number of steps for each class. The
values are rounded.

\begin{center}
\begin{tabular}{|r|r|r|r|}
\hline
$n$ & $n$ & $n \log_2 n$ & $n^2$ \\
\hline
10 & 10 & 33 & 100 \\
\hline
100 & 100 & 664 & 10,000 \\
\hline
1,000 & 1,000 & 9,966 & 1,000,000 \\
\hline
1,000,000 & 1,000,000 & about 20,000,000 & 1,000,000,000,000 \\
\hline
\end{tabular}
\end{center}

When $n$ is a million, a quadratic algorithm does a trillion steps, while a
linearithmic algorithm does about twenty million. That is why Big-O matters
more than the exact count for large inputs.

\begin{Exercise}[label=sorting1, title=Trace selection sort]
Trace the first pass of selection sort on \texttt{[4, 1, 3]}. What is the list
after the first pass? What is the final sorted list?
\end{Exercise}
\begin{Answer}[ref=sorting1]
The first pass finds that 1 is the smallest item, and swaps it with the 4. The
list becomes \texttt{[1, 4, 3]}. The second pass swaps 4 and 3, giving the
final list \texttt{[1, 3, 4]}.
\end{Answer}

\begin{Exercise}[label=sorting2, title=Count the comparisons]
How many comparisons does selection sort make on a list of 5 items, no matter
what order the items are in?
\end{Exercise}
\begin{Answer}[ref=sorting2]
$\frac{5 \cdot 4}{2} = 10$ comparisons.
\end{Answer}

\begin{Exercise}[label=sorting3, title=Prove a Big-O bound]
Show that $f(n) = 3n^2 + 100n + 5$ is $O(n^2)$ by giving constants $c$ and
$n_0$.
\end{Exercise}
\begin{Answer}[ref=sorting3]
Take $c = 4$ and $n_0 = 101$. For $n \geq 101$, we have $100n + 5 \leq n^2$,
because $n^2 - 100n - 5 = n(n - 100) - 5 \geq 101 \cdot 1 - 5 > 0$. So
$3n^2 + 100n + 5 \leq 3n^2 + n^2 = 4n^2$.
\end{Answer}

\begin{Exercise}[label=sorting4, title=Which is faster?]
A program takes $n \log_2 n$ steps and another takes $n^2$ steps. For
$n = 1{,}000{,}000$, which one is faster, and roughly how many times faster is
it?
\end{Exercise}
\begin{Answer}[ref=sorting4]
The $n \log_2 n$ program is faster. It takes about $2 \times 10^7$ steps, while
the $n^2$ program takes $10^{12}$ steps. The first is about 50,000 times faster.
\end{Answer}
